\documentclass[12pt,a4paper]{article}

\usepackage[utf8]{inputenc}
\usepackage[T1]{fontenc}
\usepackage[spanish]{babel}
\usepackage{lmodern}
\usepackage{microtype}
\usepackage{geometry}
\usepackage{setspace}
\usepackage{parskip}

\geometry{
  top=2.5cm,
  bottom=2.5cm,
  left=3cm,
  right=3cm
}

\setstretch{1.15}

\title{De lo natural y sus ambigüedades}
\author{Vicente Domínguez Arca}
\date{}

\begin{document}

\maketitle

No es intención de este autor ocultar que la siguiente reflexión tiene más papeletas de las que me gustaría para recibir fuertes críticas. Aun así, desde mi posición, siento una enorme necesidad de darle forma.

Por mucho que diferentes culturas hayan derivado hacia perspectivas menos antropocéntricas, considero que no es necesario realizar grandes esfuerzos para aceptar que el ser humano tiende, con carácter general y más allá de cualquier debate moral, a situarse en el centro de su particular \emph{Huniverso}.

Llamo aquí \emph{Huniverso} a ese fragmento del Universo al que solamente el ser humano puede acceder, bien sensorialmente, bien mediante aquellos artefactos tecnológicos producidos para extender nuestra capacidad de sensación, observación, medición o interpretación.

Nada de esto elimina, sin embargo, algo radical: somos seres humanos y formamos parte de un sistema mayor, del mismo modo que ese sistema puede formar parte de otros sistemas mayores, y así sucesivamente.

La cuestión reside entonces en dónde situamos la referencia.

Al parecer, tendemos inevitablemente a situar nuestra realidad humana en ese centro, tanto cuando pretendemos fortalecer nuestra presencia como cuando decidimos autoseñalarnos culpables.

Nos atribuimos centralidad cuando nos consideramos superiores, pero quizá también cuando nos consideramos la anomalía que ha venido a alterar aquello que denominamos naturaleza.

En ambos casos aparece una misma posibilidad: que estemos concediéndonos una excepcionalidad que no tenemos.

De ahí hemos derivado hacia debates y reflexiones que, con demasiada facilidad, dan por sentado que unas cosas son naturales y otras no lo son.

Y me gustaría plantear desde aquí una pregunta aparentemente sencilla:

¿Qué es lo no natural?

Para intentar responder debemos permanecer necesariamente dentro de aquello que observamos desde nuestro \emph{Huniverso}.

Lo que denominamos diferentes especies son, desde una perspectiva sistémica, subsistemas que separamos conceptualmente de un sistema superior.

Esa fragmentación resulta enormemente útil para comprender la realidad, pero no elimina el hecho de que cada subsistema mantiene una relación constante con todo aquello que decidimos no incluir dentro de él.

Cualquier especie que establezcamos estará, por tanto, en continua interacción con aquello que no denominamos esa especie.

Y en esa interacción, como en cualquier otra, aparecerán causas, efectos, repercusiones y contingencias.

Si se trata de una especie biológicamente activa, con carácter general necesitará obtener energía y transformar de algún modo parte de su entorno para continuar existiendo biológicamente.

Nadie parece encontrar demasiadas dificultades en llamar a esto natural.

La pregunta, llegado este punto y sin más dilaciones, es la siguiente:

¿Acaso nuestras acciones no son, lógicamente, análogas?

¿Acaso nuestras acciones no son acciones producidas por otro de los subsistemas biológicos que hemos decidido identificar dentro del sistema?

Si consideramos que no lo son, sería necesaria una argumentación razonablemente grande para explicar dónde aparece exactamente esa ruptura entre el ser humano y aquello que denominamos naturaleza.

Pero si aceptamos que el ser humano forma parte radicalmente de ella, entonces todo cuanto surge de nuestras acciones continúa sucediendo dentro de ese mismo sistema.

Aquí aparece una de las ambigüedades que más me interesan.

Solemos utilizar la palabra \emph{artificial} para aquello producido por el ser humano.

Pero ¿qué significa exactamente esa distinción?

Si pretendemos considerar artificial aquello que ha sido intencionalmente producido mediante la modificación del entorno, el criterio parece presentar problemas inmediatamente.

Pensemos, por ejemplo, en una presa construida por castores.

Existe una transformación del entorno.

Existe selección y disposición de materiales.

Existe una acción cuyo resultado modifica las condiciones anteriores del espacio.

Y existen consecuencias sobre otros subsistemas.

Sin embargo, nadie parece tener grandes dificultades para considerar aquella presa una producción natural.

Supongamos ahora una presa construida por seres humanos.

También existe transformación del entorno, selección y disposición de materiales, planificación y consecuencias sobre otros subsistemas.

La segunda será considerada artificial.

La primera, natural.

¿Dónde reside exactamente la diferencia que permite transformar una categoría en la otra?

Naturalmente existen enormes diferencias entre ambas construcciones, pero esas diferencias pueden ser observadas, descritas y parametrizadas sin necesidad de extraer una de ellas del ámbito de lo natural.

Si la última explicación consiste simplemente en que una ha sido producida por castores y la otra por seres humanos, entonces no habremos descubierto una frontera entre lo natural y lo artificial.

La habremos construido alrededor de nuestra propia especie.

Y esto resulta paradójico.

Pretendiendo clasificar la naturaleza, habremos vuelto a colocar al ser humano en una posición excepcional.

En este sentido lógico, tampoco existe una diferencia categorial entre, digamos, el guano y el plástico por el mero hecho de que el primero sea producido por unas especies y el segundo por la nuestra.

Esto no significa, obviamente, que sean equivalentes.

Cambian sus composiciones, sus procesos de producción, su persistencia, sus funciones y, sobre todo, sus consecuencias.

Pero esas diferencias aparecen cuando observamos, parametrizamos y comparamos aquello que ambos producen dentro del sistema.

No parece necesario afirmar que uno pertenece a la naturaleza mientras el otro procede de algo exterior a ella.

El plástico puede ser enormemente problemático.

Puede permanecer durante largos periodos de tiempo, acumularse, afectar a otros organismos y alterar ecosistemas.

Pero ninguna de esas propiedades convierte al ser humano en una entidad exterior al sistema que produjo ese plástico.

Somos una especie produciendo transformaciones dentro del mismo sistema en el que producen transformaciones las demás especies que observamos.

Y esto es importante, quizá demasiado, y a menudo pasa desapercibido.

Muchas veces nos afanamos sobremanera en parametrizar y observar el impacto de nuestros productos sobre otros subsistemas que identificamos: otras especies, ecosistemas, estructuras biológicas o no biológicas.

Sin embargo, olvidamos con cierta facilidad algo que considero todavía más elemental.

Muchas de nuestras acciones pueden poner en tensión las condiciones futuras de existencia de nuestra propia especie.

Esta clave no es trivial.

Si aceptamos que formamos parte del sistema sobre el que actuamos, nuestras acciones no pueden comprenderse únicamente como intervenciones exteriores sobre una naturaleza separada de nosotros.

Somos naturaleza actuando sobre naturaleza.

Un subsistema transformando el sistema del que depende.

Y, por tanto, quizá la pregunta central no deba ser si aquello que producimos es natural o artificial.

Quizá la pregunta realmente importante sea qué consecuencias tiene aquello que producimos.

Aquí resulta necesario realizar una precisión.

No se trata de considerar al ser humano \emph{la especie}.

Se trata de reconocernos como \emph{especie}.

Ni más ni menos.

Una especie entre otras especies que identificamos dentro de nuestro \emph{Huniverso}, sometida, como ellas, a determinadas condiciones de existencia y dependiente de múltiples relaciones con aquello que la rodea.

Desde esta posición, procurar nuestra supervivencia no implica afirmar que nuestra especie posea un valor superior.

Implica reconocer aquello que somos.

Las especies que observamos desarrollan estrategias, comportamientos y relaciones compatibles, al menos durante algún tiempo, con su propia continuidad.

Nuestra particularidad no reside necesariamente en perseguir esa continuidad, sino quizá en la extraordinaria capacidad que hemos desarrollado para anticipar determinadas consecuencias de nuestras propias acciones.

Y esa capacidad modifica nuestra posición.

Podemos observar algunos procesos.

Podemos modelarlos.

Podemos anticipar determinadas consecuencias.

Y, al menos parcialmente, podemos modificar nuestras acciones.

Si sabemos, por ejemplo, que determinadas emisiones aumentan la concentración de gases de efecto invernadero, que ese aumento contribuye a modificar las temperaturas y que esas modificaciones pueden alterar seriamente las condiciones dentro de las cuales nuestra propia especie ha construido sus sociedades, entonces disminuir esas emisiones no necesita justificarse exclusivamente mediante una obligación moral hacia algo exterior llamado naturaleza.

Existe una razón mucho más inmediata.

Nuestra propia continuidad.

No porque seamos más importantes que las demás especies.

Sino porque somos esta especie.

Del mismo modo que no consideramos extraño que cualquier organismo actúe de manera compatible con su supervivencia, tampoco debería resultar extraño que nosotros intentemos garantizar las condiciones compatibles con la nuestra.

La diferencia está en que somos capaces, al menos hasta cierto punto, de comprender algunas de las consecuencias futuras de aquello que hacemos.

Y esa capacidad introduce responsabilidad.

Desde esta perspectiva, lo que denominamos cuidado del entorno no necesita fundamentarse únicamente en una cuestión de respeto hacia algo exterior a nosotros.

Puede comprenderse también como una obligación radical de nuestra especie consigo misma.

No porque el entorno sea irrelevante.

Precisamente porque nunca hemos estado separados de él.

No cuidamos algo completamente exterior.

Actuamos sobre el sistema del que dependemos.

No somos ni la coronación de la naturaleza ni una anomalía introducida dentro de ella.

Somos una de sus expresiones.

Una especie.

Y quizá reconocerlo exija una suerte de catarsis.

Abandonar la idea de que somos superiores, pero abandonar también la idea aparentemente contraria de que somos una presencia exterior que ha venido a corromper una naturaleza que existiría limpia de nosotros.

Ambas imágenes pueden esconder la misma excepcionalidad.

Una nos coloca por encima.

La otra nos coloca fuera.

Quizá debamos asumir simplemente que estamos dentro.

Sucede, además, que nos encontramos inmersos en sociedades en las que determinadas lógicas económicas y políticas tienden a representar al individuo como una entidad aislada, autónoma y relativamente desacoplada de las demás.

Esta representación puede resultar útil para determinados modelos, pero entra en tensión con una realidad difícilmente discutible: ningún ser humano existe de manera independiente.

Somos nodos interconectados.

Somos tejido social.

Somos parte de estructuras biológicas, materiales y ecológicas mucho mayores.

Dependemos de otros seres humanos, de otras especies, de flujos energéticos, de condiciones materiales y de sistemas cuyo funcionamiento completo ni siquiera conocemos.

Nuestra interconexión no constituye una preferencia moral.

Es una condición de existencia.

Desde aquí, este autor no considera necesario fundamentar el cuidado de nuestro entorno exclusivamente en el respeto que debamos sentir por aquello que denominamos otras especies.

No niego que ese respeto pueda existir.

Ni niego que puedan construirse argumentos sólidos que atribuyan valor propio a otros seres vivos, especies o ecosistemas.

Simplemente considero que no resulta necesario resolver previamente esas discusiones para aceptar nuestra responsabilidad.

Podemos discutir acerca del valor intrínseco de un bosque.

Podemos discutir acerca de los derechos de otras especies.

Podemos discutir acerca de la relación moral que deberíamos mantener con un ecosistema.

Pero incluso aunque nunca llegáramos a un acuerdo sobre ninguna de esas cuestiones, seguiría existiendo una constatación anterior:

dependemos de ellos.

Nuestras preocupaciones sobre nuestras acciones pueden fundamentarse, por tanto, en cómo esas acciones afectan a nuestra sociedad y, finalmente, a nuestra continuidad como especie.

Controlar determinadas transformaciones debería dirigirse, entre otras cosas y de manera fundamental, a garantizar nuestra permanencia como subsistema dentro del \emph{Huniverso}.

Para ello es necesario interiorizar que somos nodos interconectados, que somos tejido social y que nuestra existencia futura depende de relaciones que exceden radicalmente a cada individuo.

Conservar las condiciones que permiten nuestra continuidad no debería entenderse únicamente como una preferencia.

Es una responsabilidad de nuestra especie para consigo misma.

Esto no convierte al ser humano nuevamente en el centro.

Precisamente lo contrario.

Supone renunciar a necesitar ser el centro para justificar nuestra supervivencia.

No necesitamos ser la especie más importante para querer continuar existiendo.

No necesitamos afirmar que las demás especies existen para nosotros.

No necesitamos situarnos fuera del sistema.

Basta con reconocernos dentro.

La alternativa consiste muchas veces en perdernos en innumerables disputas acerca de cómo valorar cada una de nuestras acciones sobre el entorno, construir argumentos y contraargumentos sobre el valor moral de otras especies y terminar trasladando la discusión al terreno de posiciones que difícilmente encontrarán un punto común.

\end{document}